\section{Introduction}
\label{sec:intro}

The throughput of a GPU workload is typically dominated by a small number of hardware bottleneck
classes: memory bandwidth saturation, Tensor Core under-utilization, register-pressure-limited
occupancy, and wave starvation from insufficient parallelism. The compute-bound vs.\ memory-bound
taxonomy that Williams \etal~\cite{williams2009roofline} formalized for multicore architectures
remains difficult to apply at the operator level in PyTorch, not because the hardware data is
absent, but because no tooling bridges it to the framework abstraction.

\paragraph{The attribution gap.}
Framework profilers (\texttt{torch.profiler}, \nsys)
expose operator timing without hardware counters; \ncu exposes hardware
counters without operator attribution. Bridging the two requires matching records across tools
that run in separate passes, use incompatible clock domains (CUPTI GPU timestamps vs.\ \ncu
injection-mode timestamps), and produce outputs with no shared stable key. No existing tool was designed to bridge all three (\secref{sec:related}).

\paragraph{PyTorch compilation stack.}
In PyTorch's default eager execution, each \texttt{aten::} operator dispatches individually
through the operator dispatcher at runtime. Under TorchInductor, operators
for which a hand-tuned vendor library implementation exists---such as cuBLAS for matrix
multiplication or xFormers for attention---are emitted as extern calls that
preserve this per-operator dispatch. Operators without such external counterparts---elementwise
operations, reductions, normalization, and layout transforms---are fused in groups into
Triton~\cite{triton2019} kernels, JIT-compiled to CUDA, where each kernel may represent
multiple operators and individual \texttt{aten::} dispatch does not occur at runtime.

\paragraph{Our approach.}
We observe that reliable attribution does not require timestamp matching.  Because \ncu replays
kernels in the same execution order they appear in the \nsys trace, a kernel's
\texttt{(name, invocation\_index)} pair is a stable, clock-immune identifier.  \sys combines
three complementary attribution paths: NVTX temporal enclosure for \texttt{aten::}-dispatched
kernels, invocation-order matching for \ncu counter assignment, and Inductor debug-artifact
parsing for Triton-fused kernels with no NVTX bracket.  \figref{fig:overview} shows the
complete end-to-end pipeline.

\begin{figure}[H]
\centering
\begin{tikzpicture}[
  font=\small,
  box/.style={rounded corners=3pt, draw, align=center, minimum height=1.8em, inner sep=5pt},
  toolbox/.style={box, fill=blue!12, draw=blue!50!black},
  attrbox/.style={box, fill=orange!15, draw=orange!70!black},
  outbox/.style={box, fill=green!12, draw=green!50!black},
  pgobox/.style={box, fill=red!10, draw=red!60!black},
  arr/.style={-{Stealth[length=5pt]}, thick},
  darr/.style={-{Stealth[length=5pt]}, thick, dashed, gray},
  lbl/.style={font=\scriptsize\itshape, text=gray!70!black},
]

%% ---- Row 1: Profiling stage ----
\node[toolbox] (workload) at (0,0)   {\texttt{workload.py}};
\node[toolbox] (nsys)     at (2.7,0) {\nsys capture};
\node[toolbox] (sqlite)   at (5.5,0) {SQLite export};
\node[toolbox] (ncu)      at (8.2,0) {\ncu replay};

\draw[arr] (workload) -- (nsys);
\draw[arr] (nsys)     -- (sqlite);
\draw[arr] (sqlite)   -- (ncu);

%% ---- Row 2: Attribution engine ----
\node[attrbox] (iom)    at (8.2,-1.65) {Invocation-Order\\[-1pt]Matching};
\node[attrbox] (sit)    at (5,-1.65) {StreamIntervalTree\\[-1pt]NVTX enclosure};
\node[outbox]  (prof)   at (1.5,-1.65) {\texttt{profile.json}\\[-1pt](20 counters)};

\draw[arr] (sqlite) -- (iom);
\draw[arr] (ncu)    -- (iom);
\draw[arr] (iom)    -- (sit);
\draw[arr] (sit)    -- (prof);


%% ---- Row 3: PGO ----
\node[pgobox] (diag)    at (2.7,-3.3)  {Diagnosis\\[-1pt]Engine};
\node[pgobox] (passes)  at (5.5,-3.3)  {Optimization\\[-1pt]Patterns};
\node[pgobox] (compile) at (8.2,-3.3)  {\texttt{torch.compile}\\[-1pt]+ Inductor};
\node[outbox] (opt)     at (10.8,-3.3) {Optimized\\[-1pt]model};

\draw[arr] (prof)    -- (diag);
\draw[arr] (diag)    -- (passes);
\draw[arr] (passes)  -- (compile);
\draw[arr] (compile) -- (opt);


%% ---- Feedback loop ----
\draw[darr] (opt.south) -- ++(0,-0.45) -| (workload.south)
  node[midway, below, lbl] {closed-loop validation };

\end{tikzpicture}
\caption{\textbf{\sys system overview.}  Two profiling stages combine to produce \texttt{profile.json}: \nsys
  captures the kernel timeline and \ncu captures hardware counters. The attribution engine bridges the two tools via invocation-order matching and per-stream NVTX interval trees, constituting the core contribution of \sys.  The lowest ``Optimization'' row shows a downstream application of the attributed profile: hardware-evidence-motivated
  FX graph passes validated via re-profiling; it is not part of the attribution pipeline itself.}
\label{fig:overview}
\end{figure}

\paragraph{Contribution.}
We present \textbf{\sys}, a hardware attribution pipeline that automatically links \ncu hardware
counters to PyTorch \texttt{aten::} operators via three complementary paths: NVTX temporal
enclosure, invocation-order matching, and Inductor fusion-map enrichment. Clock-domain joins are
avoided entirely; unattributed kernels are surfaced explicitly, never dropped. The pipeline
collects a curated 20-counter metric set spanning all roofline bottleneck axes via
application-mode \ncu replay, with architecture-specific fallbacks for Ampere, Hopper, and
Blackwell, and supports layer deduplication to reduce replay time by
$O(\text{num\_duplicate\_layers})$ for models with repeated structure.
